\title{Mapping AI into Production: A Field Experiment on Firm Performance}
\begin{document}
\maketitle

\begin{abstract}
Working Paper 2026/20/STR Working Paper is the author’s intellectual property. It is intended as a means to promote research to interested readers. Its content should not be copied or hosted on any server without written permission from publications.fb@insead.edu Find more INSEAD papers at https://www.insead.edu/faculty-research/research Copyright © 2026 INSEAD
\end{abstract}

\section{Recovered Text}
Working Paper
2026/20/STR
Working Paper is the author’s intellectual property. It is intended as a means to promote research to interested
readers. Its content should not be copied or hosted on any server without written permission from
publications.fb@insead.edu
Find more INSEAD papers at https://www.insead.edu/faculty-research/research
Copyright © 2026 INSEAD
Mapping AI into Production:
A Field Experiment on Firm Performance
Hyunjin Kim
INSEAD, h@insead.edu
Dahyeon Kim
INSEAD, dahyeon.kim@insead.edu
Rembrand Koning
Harvard Business School, rem@hbs.edu
March 30, 2026
AI can deliver productivity gains on individual tasks, yet evidence on whether these gains aggregate to firm
performance remains limited. We study a central friction in AI adoption, which we call the mapping problem:
discovering where and how AI creates value within a firm’s production process. Across 515 high-growth startups,
we run a field experiment in which treated firms receive information about how other firms have reorganized
production around AI, prompting them to search for use cases across a broader set of firm functions. We find that
treated firms discover more AI use cases, a 44\% increase, concentrated in product development and strategy.
These changes result in economically meaningful performance gains. Treated firms complete 12\% more tasks,
are 18\% more likely to acquire paying customers, and generate 1.9x higher revenue. Revenue and investment
gains are largest at the 90th percentile and above, consistent with AI expanding the upper range of what firms
achieve rather than modestly improving marginal ventures. Despite faster growth, treated firms do not scale inputs
proportionally. Their demand for external capital investment falls by 39.5\% relative to the control group, while their
demand for labor remains unchanged. These results provide causal evidence that AI improves firm performance
and productivity even at its current capabilities, and that discovering where and how to deploy AI is a key bottleneck
in realizing the gains from this technology.
Key words: Artificial Intelligence; Generative AI; Firms; Strategy; Entrepreneurship
Electronic copy available at: https://ssrn.com/abstract=6513481
The authors gratefully acknowledge financial support from INSEAD. The authors thank Aasiya Ahmed, Sadig
Huseynzade, Ashish Kulkarni, Simonne Pinto, Bingrui Zhong for excellent research assistance. We thank Ajay
Agrawal, Rebecca Henderson, Phanish Puranam, Michelle Rogan, and Scott Stern, along with seminar
participants at MIT Sloan, ESSEC, LBS, UCL, University of Oxford, the Singapore Initiative on Digital Economics
Workshop, and the INSEAD brownbag for their helpful feedback. The paper was formerly entitled “Reimagining
the Firm with AI: Experimental Evidence from Startups.” This experiment was approved by the INSEAD Institutional
Review Board and preregistered on the AEA RCT Registry under trial number AEARCTR-0016746.
1
Introduction
A central promise of AI rests in the idea that it will make firms more productive and drive economic
growth. The breadth of activities that AI can reshape within the firm—from operations (Brynjolf-
sson, Li, and Raymond, 2023) to strategic decisions (Csaszar, Ketkar, and Kim, 2024)—makes this
promise especially compelling. While much task-level evidence shows productivity gains across a
wide range of firm activities (Noy and Zhang, 2023; Dell’Acqua et al., 2023), these gains have largely
yet to aggregate to firm-level performance, echoing a pattern from earlier general-purpose technolo-
gies (David, 1990; Bresnahan et al., 1996). Theoretical models of the impact of AI on organizations
highlight the challenges of improving firm-level outcomes (Brynjolfsson, Rock, and Syverson, 2021;
Gans and Goldfarb, 2026). Compounding these theoretical considerations, the rapid adoption of AI
and the wide variation in tools used across firms has made it empirically difficult to isolate how AI
affects firm performance (Seamans and Raj, 2018).
In this paper, we estimate the causal effect of AI use on firm performance by leveraging a
key friction, which we call the mapping problem: the challenge of discovering where and how AI
creates value within a firm’s production process. Prior general-purpose technologies reshaped con-
crete processes—electricity demanded new factory layouts, the internet required new distribution
channels—and a large literature has studied the long-run complementary investments needed to
restructure production around them (David, 1990; Milgrom and Roberts, 1990; Bresnahan et al.,
1996). With AI, the space of possible applications is especially vast and flexible: the same technol-
ogy can assist with writing an email or fundamentally change how a product is developed, yet even
closely related tasks can differ sharply in how well AI performs (Dell’Acqua et al., 2023; Vafa, Ram-
bachan, and Mullainathan, 2024; Otis et al., 2024). As a result, a central friction is the discovery
of how to reorganize, not just the investments in reorganization itself. Discovering where and how
AI creates value is fundamentally a search problem, and a large strategy literature suggests that
search is often local (Levinthal and March, 1993; Stuart and Podolny, 1996; Kim and Wu, 2025),
especially when confronted with jagged, combinatorial, and vast landscapes like those introduced
by AI (Fleming, 2001). This suggests that firms may default to obvious local applications, while
higher-value uses remain undiscovered.
We leverage this mapping friction to design a treatment that induces exogenous variation in
firms’ ability to discover how to use AI in production. In a 3-month accelerator with 515 high-growth
startups from across the globe, we randomize the information provided through the accelerator’s
2
workshops: treated firms received case studies showing how other firms have reorganized production
around AI. These case studies included discussions of how AI-native firms have reorganized their
workflows, teams, business models, and financing to maximize the benefits of using AI (Koning,
Keller, and Kim, 2025; Koning, Kim, and Keller, 2025).
Beyond these workshops, treated and
control ventures all received standard accelerator support (Cohen et al., 2019), including general
entrepreneurship case studies, technical AI training, API credits, mentorship, and investor expo-
sure. As a result, our design lets us both estimate the causal impact of AI on firm-level outcomes
and isolate the role of mapping frictions in explaining why AI’s firm-level impact has lagged its
capabilities.
Our treatment leads to changes in how firms use AI, which in turn produce economically mean-
ingful performance gains.
Treated firms discover 2.7 additional AI use cases (a 44\% increase),
which span a broader set of activities across the firm and are especially concentrated in product
development and strategy-related domains. These changes in AI use lead to measurable gains in
performance: treated firms complete 12\% more tasks, are 11 percentage points (18\%) more likely
to acquire paying customers, and ultimately generate 1.9x higher revenues compared to control
firms. Revenue and investment gains are largest at the 90th percentile and above, consistent with
AI helping expand the upper range of venture performance by unlocking bottlenecks. These gains
are broad-based: we find no significant differential effects by baseline firm performance or founder
technical background, consistent with the mapping problem binding across a wide range of firms.
Despite faster growth, treated firms do not scale inputs proportionally: their demand for external
investment falls by just over \$220,000 relative to the control group, while their demand for labor
remains unchanged, consistent with the idea that AI enables firms to produce more output with the
same resources.
Our design also allows us to rule out a range of non-AI related explanations and to estimate the
returns to AI use within the firm. Attrition from the accelerator was low (1.6\%, eight ventures) and
balanced across treatment and control, suggesting similar engagement across groups. Analysis of
participant feedback shows that control and treated founders expressed comparable sentiment about
the accelerator experience, further reducing concerns about motivation-driven effects. Moreover,
differences between control and treated firms are driven by changes in AI use cases, making it unlikely
that the results reflect non-AI channels. Instrumenting AI use cases with treatment assignment
suggests that each additional AI use case prompted by treatment leads to 0.85 more completed tasks
and approximately 26\% higher revenue. These are large effects, suggesting that AI is fundamentally
3
reshaping how ventures scale when they can map it across their production process.
Our results make three contributions. First, we provide causal evidence that AI’s task-level pro-
ductivity gains can aggregate to firm-level performance — and identify a key friction that explains
why they often do not. Despite widespread evidence documenting productivity gains at the task
level (Brynjolfsson, Li, and Raymond, 2023; Noy and Zhang, 2023; Dell’Acqua et al., 2023), the
question of whether these aggregate to firm-level outcomes has remained open—theoretically, be-
cause complementarities and intangible investments can offset or delay returns (Gans and Goldfarb,
2026; Brynjolfsson, Rock, and Syverson, 2021), and empirically, because existing firm-level evidence
is largely correlational (Babina et al., 2024; McElheran et al., 2025; Kim and Koning, 2026). We
show that AI’s benefits do aggregate to the firm level, with gains that are broad-based across a wide
range of firms.
Second, we show that the binding constraint on AI’s value at the firm-level is not access to the
technology but search over where to deploy it. We call this the mapping problem. Most existing
evidence on the impact of AI comes from variation in whether firms adopt AI at all. We instead
generate variation in how firms map AI into their production processes. Our treatment reveals that
firms tend to search too narrowly: without examples or structured frameworks that broaden the
space they consider (Kim and Wu, 2025), firms default to familiar applications and miss higher-
value uses embedded in their production processes, consistent with a large literature in strategy
on boundedly rational search (Gavetti and Levinthal, 2000; Levinthal and March, 1993) and with
evidence that firms routinely fail to act on accessible but unattended information (Kim, 2025).
This idea also sheds light on complementarities in the adoption of general-purpose technologies
(Bresnahan et al., 1996; Brynjolfsson, Rock, and Syverson, 2021). Much of this literature emphasizes
implementation costs and lags as frictions that delay the gains from new technologies. We highlight
an earlier bottleneck: before firms can invest in redesigning their processes, they must first discover
where the complementarities lie.
Third, we show that AI is reshaping how firms are scaled.
Treated ventures achieve faster
growth without proportional increases in labor or capital, consistent with a reduction in the costs of
experimentation and scaling seen in earlier technological waves (Ewens, Nanda, and Rhodes-Kropf,
2018; Dushnitsky and Stroube, 2021).
But our upper-tail results are harder to explain by cost
reduction alone, and are instead consistent with complementarities across activities within the firm
amplifying the returns to AI (Kremer, 1993; Gans and Goldfarb, 2026). These patterns suggest
that AI is not just lowering the costs of entry or modestly improving middling ideas (Reimers and
4
Waldfogel, 2026), but expanding the frontier of what firms can be built.
2
The Mapping Problem
A growing body of task-level experiments shows that AI produces meaningful productivity gains
across a wide range of work settings (Brynjolfsson, Li, and Raymond, 2023; Noy and Zhang, 2023;
Dell’Acqua et al., 2023).
Yet, aggregating these gains to the firm level is far from automatic.
A firm’s output depends on the joint quality of many interconnected, complementary activities
(Kremer, 1993), and a large literature on general-purpose technologies shows that realizing their
benefits requires reorganizing production around them — costly investments in process redesign,
restructuring, and new skills that take time to manifest (David, 1990; Milgrom and Roberts, 1990;
Bresnahan et al., 1996). Without these investments, a task-level improvement may be bottlenecked
by unchanged processes elsewhere in the firm, producing limited or no gains in overall performance
(Brynjolfsson, Rock, and Syverson, 2021; Gans and Goldfarb, 2026). Realizing the returns to AI
thus requires not just adopting the technology but reorganizing production around it.
This complementary structure implies that before firms make these investments to benefit from
AI, managers must solve a search problem: they must discover which activities within their firm’s
production process AI can improve and how to adjust the rest of the firm. We call this the mapping
problem.
The mapping problem is difficult for three reasons.
First, AI’s capabilities are uneven and
hard to predict. Even closely related tasks can differ sharply in how well AI performs them, and
people—including experts—systematically misjudge where AI will help (Dell’Acqua et al., 2023;
Vafa, Rambachan, and Mullainathan, 2024; Otis et al., 2024). This makes it hard for managers to
assess, ex ante, which parts of their production process would benefit from AI. Second, the search
space is vast. Within any given firm, AI could potentially be applied to dozens of activities. A large
literature in strategy shows that managers search locally in the neighborhood of what they already
know (Gavetti and Levinthal, 2000; Levinthal and March, 1993), absent scaffolding like analogical
reasoning and structured frameworks (Gavetti, Levinthal, and Rivkin, 2005; Kim and Wu, 2025).
Applied to AI, this means firms default to obvious applications—a chatbot for customer support,
an LLM to draft emails—while the highest-value applications that require rethinking how work is
organized are likely to go undiscovered. Third, complementarities compound the difficulty. When
activities within the firm are tightly coupled, the value of applying AI to one activity depends on
5
whether adjacent activities also adjust (Siggelkow, 2002; Bresnahan et al., 1996). This means that
partial automation can preserve bottlenecks rather than relieve them (Demirer et al., 2025).
Consider a firm that uses AI to accelerate product development—generating a product require-
ments document or writing code faster. These are the most common and straightforward applica-
tions. But to find measurable gains, the firm must also discover that AI can be used to rapidly
prototype and test features with real users, compressing an often months-long feedback loop into
days. And if the firm does not also rethink how it identifies what customers actually want, it simply
builds the wrong product faster. The bottleneck has moved, not disappeared. When a firm solves
this discovery problem—rethinking its full production process and finding where AI can improve
not just one activity but unlock value in adjacent ones—the gains compound (Kremer, 1993; Gans
and Goldfarb, 2026). When it does not, task-level productivity gains dissipate before they reach
the firm’s bottom line.
Importantly, the mapping problem is a friction in discovery, not in access.
Two firms with
identical tools, training, and budgets can realize very different returns if one searches more broadly
across its production process for where the technology creates value. This distinguishes the mapping
problem from the access and skill frictions that dominate both the current policy conversation and
the existing empirical literature on technology adoption. Providing firms with a new technology
and teaching them how to use it may be insufficient; what matters is whether they can discover
where in their production process to deploy it and how to organize the rest of the firm to benefit
from it. While this friction is general—it arises whenever a technology’s potential applications are
vast, uneven, and hard to predict ex ante—it is especially acute for AI, where the same model can
draft an email or redesign a product and even experts cannot reliably distinguish which use will
create more value ex ante.
We study the mapping problem in the context of early-stage startups in an accelerator setting,
which offers an unusually clean context to isolate this friction. We describe this setting next.
3
Empirical Context
We designed and ran our field experiment within the AI Founder Sprint, a three-month global,
virtual startup accelerator at INSEAD designed to support early-stage ventures in building their
companies with AI. The accelerator was structured around six components—four shared equally
across all firms, and two whose content varied between treatment and control. Table 1 provides an
6
overview of the program and timeline.
Four components were shared by all firms. First, resources: participants received API credits,
access to frontier AI models, and onboarding sessions from partners including Google Cloud, Ope-
nAI, NVIDIA, and Manus AI, totaling approximately \$25,000 in-kind per firm. Second, technical
training: participants attended weekly three-hour hands-on sessions delivered by the MIT/Harvard
Sundai AI Club covering rapid prototyping, context engineering, RAG, agentic workflows, and vibe-
coding. Third, demo days and pitch practice: all firms had opportunities to pitch to leading venture
capitalists from across the globe during the Sprint’s VC Week, and were equally eligible to compete
for \$100,000 USD in non-dilutive seed funding at Demo Days in Singapore and Abu Dhabi. Fourth,
all firms submitted weekly progress reports describing their activities, financial progress, and new
AI use cases. These reports helped ventures stay on track and helped the accelerator team identify
and respond to each firm’s evolving needs. They also provided a rich source of data on venture
traction and activities, which we describe in detail below.
Starting in the third week of the accelerator, two components varied between the treatment and
control groups. First, workshops: all firms were required to attend a weekly 60-minute faculty-led
session on Zoom that delivered a case-based entrepreneurship curriculum focused on standard best
practices for making progress on early-stage venture ideas. Alongside these required workshops,
all firms were invited to six optional 90-minute masterclass workshops with leading entrepreneurs,
investors, and AI researchers.1 As discussed below, starting in the third week, treated entrepreneurs
received different information in their required workshops and were eligible for additional optional
masterclass workshops.
The final component, peer learning and office hours, was available to all firms. All ventures had
required weekly peer group meetings comprising four to eight founders with a venture guide, where
they shared progress, challenges, and plans. They could all also sign up for one-on-one or small-group
office hours with experienced mentors and accelerator instructors. These sessions were intended to
help translate information provided in the workshops into firm-specific action (Cohen and Koning,
2024). Since these sessions enabled participants to discuss and work together to implement the
workshops’ content, and since the workshop content varied across treatment and control, the peer
learning sessions naturally reinforced any differences between the two groups.
This setting provided an ideal context to study frictions in mapping AI to firm activities and
1Speakers included Ethan Mollick (Wharton), Grant Lee (Gamma), Hian Goh (Openspace Capital), Gaurav Jain
(Afore Capital), Bill Tai (Angel Investor), Claudia Zeisberger (INSEAD), and Chris Yeh (Blitzscaling).
7
processes. Early-stage ventures face relatively low organizational inertia, so difficulty in identify-
ing and integrating high-value AI use cases in this context is informative—and, if anything, likely
understates the magnitude of these frictions in larger firms. The accelerator context also provided
unusually rich measurement: we observe a large sample of ventures and can track, week by week,
how they apply AI across firm functions, what tasks they complete, and multiple dimensions of trac-
tion and performance. Finally, the program’s global, online structure supports clean experimental
variation with limited opportunities for spillover effects.
3.1
Sample selection
Ventures were selected through an open application process: any promising pre-seed venture could
apply, regardless of affiliation with INSEAD. Applications collected information on founder char-
acteristics, venture traction, and current AI usage. From the selected ventures, we restricted the
experimental sample to those who provided research consent and completed baseline data collection
in the first and second weeks of the accelerator, yielding a sample of 515 firms. Because treatment
did not begin until the third week, all sample restrictions are based on pre-treatment data.
The median firm was founded in 2024 with a team of four. The sample spanned the globe, as
shown in the map displayed in Figure A1: 33.8\% of ventures were in the Middle East and Africa,
23.6\% in Asia-Pacific, 22.3\% in Europe, and 20.2\% in the Americas. Many firms had achieved initial
traction at the time of application: 52.6\% had launched a product, 55.3\% had acquired customers,
34.6\% had generated revenue, and 35.3\% had raised external investment. Among revenue-generating
firms, median monthly revenue was \$1,000; among funded ventures, the median investment raised
was \$10,000.
3.2
Experimental Design
3.2.1
Randomization
We implemented stratified randomization to assign enrolled firms to control and treatment condi-
tions, stratifying on three pre-treatment characteristics: geography, traction score, and baseline AI
use. Geography was defined based on firm location and categorized into four regions: Middle East
and Africa, the Americas, Asia-Pacific, and Europe. Traction score was defined as the sum of three
binary indicators capturing whether the firm had launched a product, generated revenue, and raised
external investment, yielding values ranging from 0 to 3. Baseline AI use was defined as a binary
8
indicator based on the number of self-reported AI use cases at baseline: firms reporting 0-2 use
cases (below the sample median) were classified as low AI use, and those reporting three or more
use cases (at or above the median) were classified as high AI use.
These variables produced 32 strata, which had a minimum of six firms per stratum. Within
each stratum, firms were randomly assigned to treatment or control, resulting in 255 treatment
firms and 260 control firms. As shown in Table 2, baseline characteristics are well balanced across
both groups.
3.2.2
Treatment Intervention: Mapping AI into Production
Identifying where AI will perform well is difficult even for experts (Dell’Acqua et al., 2023; Vafa,
Rambachan, and Mullainathan, 2024), and discovering how to change firm-specific production pro-
cesses to benefit from AI is harder still (Brynjolfsson, Rock, and Syverson, 2021; Kim et al., 2024;
Otis et al., 2024).
Our intervention was designed to overcome these challenges by providing treated firms with
information on how other firms have reorganized around AI. To deliver this information, we leveraged
the accelerator’s workshops as a setting where we could shift what treated firms learned relative
to control. During the weekly workshop sessions, treated firms received case studies focused on
AI-native firms illustrating how they are reorganizing around this technology.2 Crucially, these case
studies went beyond which AI tools to use or the specific examples in order to help ventures abstract
and generalize to their own firm (Ruef, 2002), building on evidence that frameworks can expand
the set of alternatives visible to decision-makers (Kim and Wu, 2025). They also illustrated where
AI can be applied and how the rest of the firm must change to make such applications effective.
These discussions raised questions like whether AI demands different kinds of talent, how it changes
the cadence of product development and operational processes, and whether it opens the door to
fundamentally different business models. While these conversations were wide-ranging—they had to
be, because best practices are still being established in this space—they shared a common thread of
helping ventures think about how AI reshapes not just a single task but the structure of production
in their firms.
To illustrate the content that treated firms received, Figure 1–Figure 4 present four diagrams
2Control founders also engaged in weekly case discussions, but these focused on general entrepreneurship topics such
as building an ideal customer profile, designing tests to validate ideas, and other common themes drawn from the
lean startup and scientific approach to entrepreneurship literatures (Eisenmann, Ries, and Dillard, 2012; Camuffo
et al., 2020).
9
corresponding to four core case discussions. Drawing on recent work that explores the impact of
AI on chains of tasks within production processes (e.g., Demirer et al., 2025), each diagram uses a
simple visual structure to show how production was reorganized relative to a version without AI. In
each diagram, tasks are arranged in sequence with the responsible person (or AI) labeled beneath;
Panel A shows the production sequence before AI and Panel B how it changes with AI.
The first two cases focus on product development, a central activity for any startup. Figure 1
presents Gamma, an AI-native presentations company that has grown to over \$50M in annual re-
curring revenue with roughly 50 employees (Koning, Keller, and Kim, 2025). As Panel A illustrates,
conventional product development requires a chain of specialists, each handling one step in cycles
that can take months. With AI (Panel B), the chain compresses: AI detects usage patterns and
generates product variants directly, enabling a single PM to continuously ship features that would
previously have required an entire team. The reorganization also introduces a new task, developing
AI evals, to ensure the rapidly generated variants actually improve the product. The key insight is
that the gains come not from adopting AI for a single step but from deploying it across multiple
steps and rethinking what that combination makes possible.
Figure 2 draws on a session led by Jordan Metzner of Ryz Labs on how AI changes prototyping.
As Panel A shows, conventional prototyping is a sequential commitment: choose one tech stack,
hire one engineering team, build one version often over months, then test. With AI (Panel B),
the founder writes a Product Requirements Document and feeds it into multiple AI coding tools
simultaneously, building the same idea multiple ways rather than betting on a single approach. This
creates a new bottleneck: with building no longer the constraint, the limit becomes how quickly the
founder can get customer feedback. Here too, AI helps: because prototypes update in minutes, the
founder can test the best version with a customer, build a feature live on the call, and test again.
The remaining two cases illustrate how the mapping problem extends well beyond product devel-
opment to business models and operations. Figure 3 presents FazeShift, an AI accounts receivable
startup whose co-founder Caitlin Leksana led a session on workflow automation (Koning, Kim, and
Keller, 2025). Panel A shows the conventional accounts receivable process: eight steps alternating
between software systems and humans who bridge between them. With AI (Panel B), this bridging
is replaced by a fully automated sequence. The key insight is that partial automation preserves
the bottleneck rather than relieving it (Demirer et al., 2025); the full chain of steps must be auto-
mated to convert what was previously a labor-intensive professional service into a scalable software
product.
10
Figure 4 presents a case on Ranger, illustrating how reorganizing around AI can shift the business
model itself. As Panel A shows, the conventional VC-funded path front-loads capital: the chain
begins with “Raise money” and selling does not happen until step four. Panel B reorders the entire
chain. The founder starts by selling QA testing as a service from day one, building tacit knowledge
about which steps are routine and where AI can automate. By using AI to improve their own
efficiency and that of future employees, the founder transforms the margins of the business from a
services firm to something closer to a product firm. The most visible shift is that “Raise money”
moves from the first box in Panel A to the last in Panel B, giving the venture optionality on whether
and when to raise outside capital.
Peer learning and office hours reinforced the workshop content. As described above, all firms
had required weekly peer group meetings and optional office hours with mentors and accelerator
instructors. Because the workshop content differed across treatment and control, and peer sessions
enabled participants to discuss and implement what they learned, these sessions naturally amplified
treatment-control differences.
Finally, to minimize potential spillovers, treatment and control participants were placed into
separate groups on the program’s online platform and were unable to directly message or interact
with members of the other group. During shared program components, interactive features including
the Zoom chat function were disabled to prevent cross-group communication.
These measures
limited opportunities for participants to exchange information about workshop content or access
materials outside their assigned condition.
4
Data and Empirical Specification
4.1
Data
We collected a wide range of data throughout the accelerator, including both quantitative measures
of venture development and qualitative evidence of founders’ behavior and thinking as reflected in
text data.
Baseline application materials.
Prior to the accelerator, all ventures submitted application
materials, including a description of their business, a pitch video, self-reported growth metrics,
and additional information about the team. These data provide the basis for our pre-treatment
covariates, enabling stratified randomization and the assessment of baseline balance across treatment
11
arms.
Weekly progress reports.
Our primary data source is a series of structured ‘weekly progress
reports’ submitted by each firm throughout the Sprint. Each week, ventures completed a survey
that collected data on venture progress, which reported: (1) growth and financial metrics including
the number of customers, revenues, costs, and investments; (2) their anticipated demand for capital
and labor by the Demo Days in December (over 3 months away from the start of the program); (3)
up-to-date descriptions of their ventures and whether they had pivoted; (4) any new AI use cases in
the last week; (5) specific tasks or actions they took to advance the startup that week; (6) product
demos showing the changes they made; (7) in two of the weeks, pitch decks and videos. This weekly
cadence of detailed data collection allows us to observe venture progression over time and construct
cumulative measures of activity over the post-baseline period, in addition to endline snapshots of
venture performance. Compliance was high: 98.4\% of ventures (507 of 515) completed at least one
post-baseline weekly progress report, and 81.0\% completed all eight progress reports (Table A1).
Process \& engagement data.
We also collected supplementary data on venture engagement
throughout the program. This includes recordings and transcripts of weekly guided peer group
meetings, workshop attendance records, and discussion logs from our online community forum.
These sources are used primarily for compliance checks and mechanism analyses rather than as
primary outcome measures.
4.2
Outcome Measures
Our primary measure of AI usage is the number of distinct AI use cases reported by each firm,
cumulated over the post-baseline period. Each week, founders described in their own words any new
AI use cases they had implemented or explored. All responses were independently reviewed by two
independent human coders, who assessed whether each entry constituted a valid AI application and
flagged duplicates (i.e., entries already reported by the same firm in a previous week). Disagreements
were resolved by a third independent coder.3 Full coding guidelines are provided in Appendix 7.2.1.
We also measure the pace and direction of venture activity. Each week, ventures reported the
concrete tasks they completed to advance the startup. The number of concrete tasks completed is
constructed analogously to AI use cases — cumulated over the post-baseline period, coded by two
3All coders were blind to the experiment, hypotheses, and treatment conditions.
12
independent human coders and resolved by a third, and classified as either internal (e.g., product
development, coding) or external (e.g., customer meetings, investor pitches) (Appendix 7.2.2).
We study a number of measures of venture performance. First, we construct an index of venture
progress, which represents the sum of four binary indicators: whether the venture had launched its
product, acquired customers, generated revenue, and raised investment at endline. We standardize
this sum using the full-sample mean and standard deviation so that treatment effect estimates are
interpretable in standard deviations. We also examine each of the four binary variables separately.
In addition to these progress metrics, we measure the total amount of revenue generated, team size
and composition, and investment raised for each venture by endline, our last week of data collection.
We also track whether ventures pivoted from their original idea at any point during the Sprint and
whether they were shut down at endline.
Finally, we measure ventures’ estimated demand for inputs. We measure capital demand as the
amount ventures reported they would seek to raise on Demo Day, and labor demand from founders’
responses describing roles they would like to post job postings for. Each listed role was extracted
from free-text, screened to exclude any invalid entries unrelated to hiring, and matched against a
keyword list to assign a functional category, with tech-related postings defined as engineering, AI
specialist, product, or design roles. We categorize job postings and examine differences in demand
for tech-related roles to assess whether treatment meaningfully shifts the types of roles founders are
seeking to hire.
4.3
Empirical specification
Our main specification estimates the intent-to-treat (ITT) effect of assignment to treatment:
yi = β1 Treatmenti + δs(i) + εi
(1)
where yi is the outcome of interest for firm i, Treatmenti is an indicator for treatment assign-
ment, and δs(i) are fixed effects for randomization strata defined by geography, baseline traction,
and baseline AI use. Standard errors are clustered at the firm level, the unit of randomization.
Heterogeneity Analysis.
To examine heterogeneous treatment effects, we estimate specifications
of the form:
yi = β2 Treatmenti + γ Xi + ϕ (Treatmenti × Xi) + δs(i) + εi
(2)
13
where Xi is a pre-treatment firm or founder characteristic capturing treatment effect hetero-
geneity. The coefficient ϕ captures differential treatment effects as a function of Xi. We examine
heterogeneity across pre-registered stratification dimensions such as baseline traction and baseline
AI use, and also conduct exploratory analyses on additional founder and firm characteristics such
as technical background, prior startup experience, and team size.
Instrumental Variables Specification.
Our ITT estimates capture the effect of being assigned
to treatment, but since treatment shifts the number of AI use cases firms adopt, we can go further
and estimate the returns to each additional use case. This matters for two reasons. First, it sheds
light on the mechanism: if treatment effects are mediated by changes in AI use, channels unrelated
to AI are unlikely to drive our results. Second, it provides a better estimate of AI’s returns, since
control firms also use AI—the ITT captures only the marginal difference in use cases between groups,
not the full value of AI adoption and integration. To recover the per-use-case effect, we estimate
two-stage least squares (2SLS) specifications that instrument firm-level AI usage with randomized
treatment assignment.
The first stage estimates the effect of treatment on AI use:
AIUseCasesi = π0 + π1 Treatmenti + δs(i) + νi
(3)
The second stage uses predicted AI use to estimate its effect on firm outcomes:
yi = β3
d
AIUseCasesi + δs(i) + εi
(4)
Because treatment was randomly assigned, it provides exogenous variation in AI adoption.
Assuming the exclusion restriction—that treatment affects outcomes only through its effect on AI
use—the 2SLS estimate β3 identifies a local average treatment effect (LATE), corresponding to the
effect of AI use cases for firms whose AI use was induced by the treatment. In practice, the exclusion
restriction is unlikely to hold exactly: treatment may have shifted not just the number but also the
nature of AI use cases adopted, for example toward applications involving deeper reorganization of
production processes. The IV estimates should therefore be interpreted as capturing the return to
the type of AI use cases induced by treatment, not the return to any arbitrary additional use case.
14
5
Results
5.1
Do treated firms change how they use AI?
We begin by examining whether the treatment changes how firms use AI. Each week, ventures
reported any new ways they were using AI in their firm, from using Claude Code to build an
automated report writer, to deploying Manus AI as a clinical affairs specialist, to prototyping an
AI-driven scenario engine in an HR dashboard (see Appendix Table A8 for additional examples).
These self-reported use cases, validated by independent human coders, are our primary measure of
how broadly firms map AI into their production processes.
As shown in Figure 5, treatment shifts the entire distribution of AI use cases. Panel (a) displays
the distribution of cumulative post-baseline use cases: the modal control firm reports zero or one use
case, while treated firms are shifted substantially to the right. On average, treated firms reported
8.8 cumulative AI use cases compared to 6.1 for control firms, a difference of 2.7 use cases or roughly
a 44\% increase relative to the control mean. Panel (b) shows that this gap emerges after treatment
begins in week 3 and grows steadily over the program, consistent with cumulative discovery rather
than a one-time shift. Our primary specification with strata fixed effects yields an estimate of +2.69
additional AI use cases (p < 0.01; Table 3).
Treatment also expands the breadth of AI adoption. We classify reported applications into seven
functional domains: technical infrastructure, product development, marketing, research/legal/writing,
product/strategy design, business operations, and sales (details in Appendix 7.3). Treated firms
adopt AI across 0.84 more distinct functional categories compared to the control (p < 0.01; Table 3),
suggesting that the intervention does not simply intensify existing uses but expands the range of
business functions in which firms apply AI. This breadth matters: when firm production is com-
plementary, applying AI to only one function can leave bottlenecks intact and mute performance
effects (Gans and Goldfarb, 2026).
Moreover, treated firms differ significantly in the types of use cases adopted. Figure 6 breaks out
average use cases per venture across seven functional categories. The largest and most significant
differences are in product development, where treated firms use AI to prototype and build products;
product/strategy design, where firms embed AI as a core feature of the product itself (Kim and
Koning, 2026) and use AI to inform high-level strategic decisions; and business operations, where
firms automate internal workflows like email triage and process coordination. The differences are
relatively smaller in categories like research/legal/writing and sales, where AI use looks more similar
15
across groups. The pattern is telling: treatment shifts AI usage most in areas where founders must
rethink how work is organized, not just layer a chatbot onto an existing task. Put differently, the
treatment expands both AI as a process input—where AI helps the team build faster—and AI as a
product output—where AI capabilities are delivered directly to customers. That effects concentrate
in both suggests the intervention expanded not just how firms build but what they build.
These results reflect differences in AI use cases, not in the AI tools firms use. Appendix Figure A5
displays AI tool adoption across treated and control ventures: all major tools are used broadly by
both groups, though treated firms do adopt more tools overall. The treatment gets firms to use
AI more, but as the use case results above show—and as we discuss in the sections that follow—it
also changes how firms reorganize the rest of their production processes around it. The binding
constraint is not access to the technology, since both groups use multiple AI tools, but the ability
to discover where they create value and how to reorganize the firm accordingly.
5.2
Does treatment change what tasks firms complete?
These changes in AI use lead to meaningful differences in venture activity. Each week, ventures
reported the concrete actions they took to advance their startup, from building an MVP to pitching
investors to onboarding beta testers (see Appendix 7.2.2 for coding details and examples). Aggre-
gating across the post-baseline period, treated firms completed significantly more tasks than control
firms—20.9 tasks on average compared to a control mean of 18.5 (Figure 7). Our primary specifi-
cation with strata fixed effects yields an estimate of +2.3 tasks (p < 0.05; Table 3). The treatment
effect grows over time, consistent with cumulative gains from increased AI use (Figure 7).
To understand where these gains occur, we had two independent human coders classify each
task as either internal—activities performed within the firm like developing a product, building a
landing page, or creating financial projections—or external—activities involving people outside the
firm like customer interviews, investor pitches, or sales outreach (Appendix 7.2.2). A third coder
resolved disagreements.
The treatment effect is driven almost entirely by internal tasks. Treated firms complete ap-
proximately 2.2 more internal tasks relative to a control mean of 9.8, a 22\% increase (p < 0.01),
with no statistically significant difference in external task completion (Table 3). This pattern makes
sense given our treatment design: the accelerator provided identical opportunities to meet investors
and potential customers across all firms, so we would not expect external tasks to differ. What AI
changes is how much firms can build—and the internal task results suggest they are building more.
16
Founders’ own reflections point in the same direction. One treated founder reflected: “This
mindset shift fundamentally changed how we build at [REDACTED]. I began using AI tools not as
a replacement for expertise but as a force multiplier . . . a way to perform the ‘analyst or intern’
version of a task that I can then refine further.” Another explained: “In just a few hours I was able
to produce what previously cost \$1,000 from an outsourced dev team . . . We realized we will need
even fewer staff than we initially thought, because AI allows us to do in house what we assumed
required outsourcing.”
5.3
Does treatment impact firm performance?
Measuring firm performance is challenging because early-stage ventures can make meaningful progress
without yet generating revenue, and costs are often unclear. In practice, investors assess progress
through a set of milestones, and we take a similar approach.
We examine whether firms have
launched a product, acquired customers, generated revenue, and raised investment—each capturing
a different dimension of traction (Table 4).
We find that treatment induces meaningful improvements across these milestones.
Overall,
treated firms’ Venture Progress Index increases by 0.16 standard deviations (p < 0.05).
Ap-
pendix Figure A4 shows the distribution of the underlying index, which sums the four binary
milestones (range 0–4). The histogram reveals that treatment both reduces the share of ventures
achieving zero milestones and shifts substantially more ventures to achieving all four. This suggests
that AI is not simply helping firms produce a prototype that nobody wants; rather, treated firms
are progressing along a range of milestones—from product launch through customer acquisition and
revenue generation—consistent with genuine advancement. Turning to the individual components,
treated firms are 4.7 percentage points (6\%) more likely to have launched a product (p < 0.1) and
11 percentage points (18\%) more likely to have acquired paying customers (p < 0.01). Table A2
further supports these results: treated firms are 8.9 percentage points more likely to transition
from having no product to launching one, and 14.4 percentage points more likely to transition from
having no customers to acquiring their first. Treated firms generate 1.9× higher total revenue on
average (log(x+1) coefficient 0.65, p < 0.1); conditional on earning revenue, the effect rises to 2.2×
(log coefficient 0.78, p < 0.05). Both results are robust to alternative transformations, including the
inverse hyperbolic sine and various pre-registered winsorization thresholds.(Appendix Table A3).
We observe limited average effects on investment outcomes, with no statistically significant impact
on the likelihood of raising external capital or on total investment amounts (Table 4; Appendix Ta-
17
ble A4).
Notably, these gains do not appear to be driven by increased pivoting or by ventures shutting
down and restarting with a new idea. The coefficient on pivoting is positive but statistically in-
significant (0.057, p > 0.1), and the coefficient on shutdown is slightly negative and insignificant
(−0.005, p > 0.1; Table A2).
Treatment improves firm performance on average, but we are often interested in how a technology
shifts the distribution of outcomes (Csaszar, Ketkar, and Kim, 2024)—and in particular, whether it
produces right-tail winners, especially for high-growth startups (Kerr, Nanda, and Rhodes-Kropf,
2014; Koning, Hasan, and Chatterji, 2022).
We estimate quantile regressions that allow us to
examine treatment effects at each point in the outcome distribution. As shown in Figure 8, the
revenue gains from treatment are small through most of the distribution but rise sharply in the
upper tail, with the largest effects emerging at the 90th and 95th percentiles. These quantile results
also help explain the noisiness of the average log revenue estimate: the treatment effect is driven by a
subset of firms that break out, rather than a uniform shift. Turning to investment, we find a similar
pattern: the right tail of the investment distribution is wider for treated firms (Figure 9), with the
highest-performing treated firms raising substantially more capital than their control counterparts.
Full quantile regression tables are reported in the Appendix (Tables A10–??).
These right-tail gains are consistent with the mapping problem framework introduced above:
when founders learn to apply AI across their production processes, it can alleviate binding bottle-
necks that unlock disproportionate value (Kremer, 1993; Gans and Goldfarb, 2026). The pattern
suggests that AI lifts the ceiling of what the most promising ventures can achieve, rather than
primarily lowering entry barriers or modestly improving outcomes for marginal firms.
5.4
Do treated firms think differently about the inputs they need?
A striking pattern in our data concerns not what treated firms produce, but what they consume.
Given faster growth and increased traction, one might expect treated firms to demand more external
resources. Yet we observe the opposite. Treated firms report just over \$220,000 less in capital
demand relative to control firms, a 39.5\% decrease (p < 0.05), with no corresponding increase in
labor demand (Table 5). This capital demand result is robust to alternative transformations and
winsorization levels (Appendix Table A5).
These averages are not to say that AI-using firms do not need to raise capital.
When we
examine the quantile distribution of investment raised, we find that treated firms in the upper tail
18
raise substantially more funding than their control counterparts (Figure 9). However, this conflates
two things: the firm’s need for capital and its ability to attract it, since faster-progressing ventures
may simply be more appealing to investors. Consistent with the idea that AI lowers firms’ perceived
need for capital, when we turn to the quantile distribution of capital demand (Figure 10), we see
that the largest reductions in capital ask come from the right tail, above the 60th percentile. Taken
together, our results suggest that treated firms think they can do more with less, consistent with
the idea that AI may fundamentally change what firms can produce with a given set of inputs. Full
quantile regression tables for capital demand are reported in the Appendix (Tables A12–??); these
patterns are robust to winsorizing at the 99th percentile.
5.5
Alternative explanations for the performance effects
A natural concern is that treatment effects reflect heightened motivation or effort rather than
the specific content of the intervention—a Hawthorne-style effect. We think this is unlikely. If
control firms were simply less motivated, they should be more likely to attrit. But as shown in
Appendix Table A1, weekly survey response rates are nearly identical across groups throughout
the program, with no statistically significant differences at any point; if anything, control firms
were slightly more likely to complete all surveys without missing a week (81.2\% vs. 80.8\%). We
also find no differences in sentiment about the program: Appendix Figure A3 reports sentiment
analysis of founders’ open-ended text responses using both VADER and RoBERTa classifiers, with
no meaningful differences in affect between treated and control firms.
A second concern is spillovers between treatment and control firms. We took several steps to
prevent this. Treatment and control firms were assigned to fully separate groups on the program’s
online platform with no ability to view each other’s content, and live chat was disabled during
joint sessions with outside speakers. When asked at endline about sessions they wished they could
have attended, zero control founders mentioned treatment-specific sessions.
Attendance records
confirm that no control firms attended treatment-only workshops. Additional details are provided
in Appendix 7.1.
5.6
How do treated firms map AI into their production processes?
Having ruled out these alternative explanations, what evidence do we have that firms actually
solved the mapping problem and reorganized around AI? The results above show that treated firms
use more AI, complete more tasks, perform better, and demand less capital. But what does this
19
reorganization actually look like? Appendix Table A9 presents illustrative examples drawn from
ventures’ weekly progress reports, in which they described how they were using AI and what they
accomplished each week. Several patterns emerge that connect directly to the mechanisms outlined
in our treatment vignettes (Figure 1–Figure 4). In each, the key step was not learning to use AI—
all firms had the same tools and training—but recognizing where AI could reshape how value was
created, a cognitive shift from seeing AI as a tool for existing tasks to seeing it as a way to rethink
the production process itself.
First, treated firms use AI to rapidly prototype and build products that would otherwise require
specialized talent.
A tender-matching platform prototyped an end-to-end pipeline from tender
classification through compliance checking and bid pricing, going from zero to \$40K in revenue
and four paying customers during the program—without hiring the technical talent that would
conventionally be needed to build this pipeline.
Second, we see firms reorganize internal workflows by replacing sequential, labor-intensive pro-
cesses with AI. A field services venture rebuilt its full operations chain—from an intent classifier
replacing a dispatcher to a payment-reminder engine replacing collections staff—compressing five
roles into AI modules that self-improve from operator feedback. A fintech with 2,500 existing cus-
tomers moved KYC screening, its call center, and even hiring interviews to AI, growing its customer
base by 20\% while reducing costs.
Third, these changes enable firms to operate at scale or in markets that were previously un-
economical. A maternal health venture replaced hour-long clinical intake sessions with a What-
sApp chatbot that conducts intake and classifies risk in under eight minutes, making it viable to
screen patients in low-resource settings where trained health workers are scarce. A trade quality-
control startup built a multi-modal pipeline—image forensics, compliance report drafting, location
verification—that reduced manual inspection by 70\% and grew the venture from two to forty cus-
tomers during the program. A medical transport company shifted from reactive human dispatch
to AI-driven routing by patient urgency and traffic, with predictive positioning of drivers before
hospital discharge peaks, reaching 160 rides per week with 60\% of bookings handled entirely by AI.
In each case, the gains came from founders discovering where and how AI could create value across
their production process—not from any single, obvious application.
20
5.7
Who benefits from mapping AI into production?
A natural question is whether the benefits of mapping and reorganizing around AI accrue broadly
or only to certain types of founders and firms. We examine heterogeneity along two dimensions:
baseline venture performance and founder technical background.
We first examine heterogeneity by baseline firm performance (Appendix Table A7). On our
primary outcome, the Venture Progress Index, the interaction between treatment and high traction
is small and insignificant, indicating no differential effects by baseline performance. The customer
acquisition effect is somewhat concentrated among firms that had not yet reached that milestone
(−19.6 pp interaction, p < 0.05), but otherwise the treatment helps across the traction distribution.
This contrasts to task-level experiments that have found that low performers benefit the most from
AI (Noy and Zhang, 2023; Brynjolfsson, Li, and Raymond, 2023; Dell’Acqua et al., 2023; Csaszar,
Ketkar, and Kim, 2024) and to recent work showing that only high-performing small business owners
have the judgment and skills required for AI to improve business performance (Otis et al., 2024).
We next ask whether technically trained founders benefit more.
If the intervention primar-
ily taught technical skills—how to prompt, how to code with AI—we would expect founders with
technical backgrounds to benefit less, since they likely already possess these capabilities. Instead,
we find no significant differential effects by technical background: the interaction between treat-
ment and having a technical founder is small and statistically insignificant across all outcomes
(Appendix Table A6).
Together, these results are consistent with the mapping problem being a cognitive constraint
on discovery that our treatment helps founders overcome. If the friction were technical, techni-
cally trained founders should benefit less. If the friction were related to general managerial or en-
trepreneurial skills, higher-performing firms should benefit more. That neither predicts treatment
response suggests the binding constraint is how broadly founders search across their production
processes—a friction that examples and frameworks appear to alleviate regardless of background or
starting position (Kim and Wu, 2025).
5.8
Estimating the firm-level returns to AI use
Our ITT estimates capture the effect of helping firms solve the mapping problem, but the interven-
tion also lets us estimate the returns to AI use itself. Because our treatment shifts the number of AI
use cases firms adopt, we can instrument firm-level AI usage with randomized treatment assignment
21
to recover the marginal value of each additional use case (Table 6). This approach identifies the
local average treatment effect of an additional AI use case for firms whose adoption was shifted by
the treatment—providing an estimate of AI’s returns that extends beyond the mapping friction per
se. It requires the assumption that treatment affects outcomes only through its effect on AI use.
The first stage is strong, with F-statistics exceeding 33, confirming that treatment substantially
shifts AI use. The second-stage estimates suggest that each additional AI use case increases the
number of tasks completed by 0.85 (p < 0.01), raises the Venture Progress Index by 0.056 standard
deviations (p < 0.05), increases the probability of having customers by 3.9 percentage points (p <
0.01), increases log revenue by 0.23 (p < 0.1), and reduces log capital demand by 0.17 (p < 0.1).
These estimates are consistent with the ITT results: the treatment-induced 2.7 additional use cases,
valued at 0.06 SD each, predicting a 0.15 SD increase in the Venture Progress Index—close to the
ITT estimate of 0.16 SD.
To put the magnitude in perspective, control firms already average 6.1 AI use cases and treated
firms average 8.8. If we take the per-use-case LATE at face value and extrapolate linearly, a firm
moving from zero AI use cases to the treatment group mean of 8.8 would see a 0.49 SD increase
in the Venture Progress Index and a roughly 34 percentage point increase in the probability of
acquiring paying customers.
These are obviously relatively heroic extrapolations—the LATE is
identified locally around the treatment-induced margin. But they provide a useful benchmark: the
cumulative value of mapping AI broadly across a firm’s production process appears to be substantial,
and the treatment-induced margin of 2.7 additional use cases captures only a fraction of the value
of AI.
6
Discussion
We study 515 high-growth startups in a field experiment that helped treated firms to explore how
AI could reshape their production processes. Treated firms identified more AI use cases, completed
more tasks, and were more likely to acquire paying customers and grow revenue—with the largest
gains in the upper tail, consistent with with the idea that reorganizing production around AI can
expand the upper range of what ventures can achieve. Treated firms also demanded hundreds of
thousands of dollars less in external capital, with no change in workforce size or labor demand,
consistent with a view that AI enables firms to produce more with fewer resources. These effects
are broad-based across technical backgrounds and performance levels, suggesting that the mapping
22
problem is a key constraint on firm-level value creation.
Taken together, these results offer a partial resolution to the disconnect between AI’s demon-
strated capabilities and its limited appearance in firm-level and aggregate productivity statistics.
Our result
\end{document}
